%% kapa: exam=1 round=22 year=2011 period=2 form=A subject=appraisal_law subject_ko=부동산관계법규 seq=11 number=51
%% source: Q-Net 감정평가사 기출문제 (cst00309) artlSeq=5036790 file=extracted/1008/5036790/2011년 제22회 감정평가사 1차시험 기출문제/2교시.hwp
%% conversion: kordoc 4.24.51 -> tools/gen1.py (auto); review=proofread-gpt-astra; flags=none
\begin{question}{51}
국토의 계획 및 이용에 관한 법령상 용도지역ㆍ지구ㆍ구역에 관한 설명으로 옳지 않은 것은?
\begin{choices}
  \item 용도지역ㆍ지구ㆍ구역의 변경은 도시관리계획으로 결정된다.
  \item 용도지역은 서로 중복되어 지정될 수 없으나, 용도지구는 서로 중복되어 지정될 수 있다.
  \item 상업지역은 중심상업지역, 일반상업지역, 근린상업지역과 유통상업지역으로 세분된다.
  \item 시가화조정구역내에서 시가화조정구역의 지정목적달성에 지장이 없다고 인정되는 도시계획사업의 시행에 관한 요청권은 당해 지방자치단체장에게 부여되어 있다.
  \item 특별시ㆍ광역시ㆍ시 또는 군의 도시계획조례로 용도지역별 용적률을 정함에 있어서 필요한 경우에는 당해 지방자치단체의 관할구역을 세분하여 용적률을 달리 정할 수 있다.
\end{choices}
\end{question}
\answer{4}
